\documentclass[8pt]{article}
\linespread{0.5} 
\usepackage{extsizes}
\usepackage{amsmath}
\usepackage{amssymb}
\usepackage{amsthm}
\usepackage{ mathrsfs }
\usepackage{ mathtools}
\usepackage{enumerate}
\usepackage{bbm}
\usepackage{lipsum}
\usepackage{fancyhdr}
\usepackage{tikz-cd} 
\usepackage{adjustbox} 
\usepackage{ stmaryrd }
\numberwithin{equation}{section}
\usetikzlibrary{arrows}
\newcommand{\midarrow}{\tikz \draw[-triangle 90] (0,0) -- +(.1,0);}
\tikzset{commutative diagrams/.cd,
mysymbol/.style={start anchor=center,end anchor=center,draw=none}
}
\newcommand\comsymb[2][\circlearrowleft]{%
  \arrow[mysymbol]{#2}[description]{#1}}

\newtheorem{theorem}{Theorem} [section] 
\newtheorem{proposition}{Proposition}[section] 
\newtheorem{definition}{Definition}[section] 
\newtheorem{lemma}{Lemma}[section] 
\newtheorem{notation}{Notation}[section] 
\newtheorem{remark}{Remark}[section] 
\newtheorem{corollary}{Corollary} [section] 
\newtheorem{terminology}{Terminology}[section] 
\newtheorem{fact}{Fact}[section] 
\newtheorem{example}{Example}[section] 
\newtheorem{claim}{Claim}
\newtheorem*{question}{Question} 
\newtheorem*{conjecture}{Conjecture}
\newenvironment{claimproof}[1]{\par\noindent\underline{Proof:}\space#1}{\hfill $\blacksquare$}

\DeclareMathOperator{\tmod}{mod}
\DeclareMathOperator{\triv}{triv}
\DeclareMathOperator{\ind}{ind}
\DeclareMathOperator{\straw}{s.t.}
\DeclareMathOperator{\ev}{ev}
\DeclareMathOperator{\pr}{pr}
\DeclareMathOperator{\Aut}{Aut}
\DeclareMathOperator{\Map}{Map}
\DeclareMathOperator{\Stab}{Stab}
\DeclareMathOperator{\Ind}{Ind}
\DeclareMathOperator{\GL}{GL}
\DeclareMathOperator{\SL}{SL}
\DeclareMathOperator{\SO}{SO}
\DeclareMathOperator{\Sp}{Sp}
\DeclareMathOperator{\esssup}{esssup}
\DeclareMathOperator{\abs}{abs}
\DeclareMathOperator{\rel}{rel}
\DeclareMathOperator{\diam}{diam}
\DeclareMathOperator{\rank}{rank}
\DeclareMathOperator{\Hom}{Hom}
\DeclareMathOperator{\Homk}{Hom_{k-alg}}
\DeclareMathOperator{\Ker}{Ker}
\DeclareMathOperator{\Image}{Im}
\DeclareMathOperator{\Dom}{Dom}
\DeclareMathOperator{\grad}{grad}
\DeclareMathOperator{\divergence}{div}
\DeclareMathOperator{\rk}{rk}
\DeclareMathOperator{\Span}{Span}
\DeclareMathOperator{\mSpec}{mSpec}
\DeclareMathOperator{\MaxSpec}{MaxSpec}
\DeclareMathOperator{\interior}{int}
\DeclareMathOperator{\supp}{supp}
\DeclareMathOperator{\id}{id}
\DeclareMathOperator{\sgn}{sgn}
\DeclareMathOperator{\Mat}{Mat}
\DeclareMathOperator{\PD}{PD}
\DeclareMathOperator{\ext}{ext}
\DeclareMathOperator{\vol}{vol}
\DeclareMathOperator{\inte}{int}
\DeclareMathOperator{\diverg}{div}
\DeclareMathOperator{\curl}{curl}
\DeclareMathOperator{\image}{im}
\DeclareMathOperator{\dR}{dR}
\DeclareMathOperator{\Ob}{Ob}
\DeclareMathOperator{\Mnfd}{Mnfd}
\DeclareMathOperator{\OpEmb}{OpEmb}
\DeclareMathOperator{\Spec}{Spec}
\DeclareMathOperator{\Spa}{Spa}
\DeclareMathOperator{\Lim}{Lim}
\DeclareMathOperator{\Alg}{\mathcal{A}}
\DeclareMathOperator{\glim}{glim}
\DeclareMathOperator{\con}{con}
\DeclareMathOperator{\Spv}{Spv}
\DeclareMathOperator{\clet}{c}
\DeclareMathOperator{\tat}{at}
\DeclareMathOperator{\Sheaf}{Sheaf}
\DeclareMathOperator{\Cont}{Cont}
%\newcommand*{\name}[\num_arguments][default values]{{\color{#1}\Large #2}}
\newcommand{\defeq}{\vcentcolon=}
\newcommand{\Ouv}{\mathscr{O}}
\newcommand{\norm}[1]{\Vert #1 \Vert}
\newcommand{\opNorm}[2]{\norm{#1}_{#2\to#2}}
\newcommand{\normL}[3]{\norm{#1}_{L^{#2}(#3)}}
\newcommand{\N}[0]{\mathbb{N}}
\newcommand{\m}[0]{\mathfrak{m}}
\newcommand{\R}[0]{\mathbb{R}}
\newcommand{\Z}[0]{\mathbb{Z}}
\newcommand{\C}[0]{\mathbb{C}}
\newcommand{\Q}[0]{\mathbb{Q}}
\newcommand{\F}[0]{\mathcal{F}}
\newcommand{\U}[0]{\mathcal{U}}
\newcommand{\B}[0]{\mathcal{B}}
\newcommand{\G}[0]{\mathbb{G}}
\newcommand{\bA}[0]{\mathbb{A}}
\newcommand{\bT}[0]{\mathbb{T}}
\newcommand{\cK}[0]{\mathcal{K}}
\newcommand{\cH}[0]{\mathcal{H}}
\newcommand{\Para}[0]{\mathbb{P}}
\newcommand{\M}[0]{\mathcal{M}}
\newcommand{\maniN}[0]{\mathcal{N}}
\newcommand{\cinf}[0]{\mathcal{C}^\infty}
\newcommand{\torus}[0]{\mathbb{T}}
\newcommand{\cR}[0]{\mathcal{R}}
\newcommand{\sep}[0]{\:|\:}
\newcommand{\pd}[2]{{\frac {\partial #1} {\partial #2}}}
\newcommand{\compinc}[0]{\subset \subset}
\newcommand{\sH}[0]{\mathscr{H}}
\newcommand{\Nfil}[0]{\mathcal{N}}
\newcommand{\g}[0]{\mathfrak{g}}
\newcommand{\p}[0]{\mathfrak{p}}
\newcommand{\QC}[0]{\mathfrak{Q}_C}
\newcommand{\ab}[1]{\vert #1 \vert}
\newcommand{\st}[0]{\:\straw\:}
\newcommand{\cGamma}[0]{c\Gamma}
\newcommand{\sC}[0]{\mathscr{C}}
\newcommand{\sU}[0]{\mathscr{U}}
\newcommand{\cV}[0]{\mathcal{V}}
\newcommand{\sA}[0]{\mathscr{A}}
\newcommand{\sB}[0]{\mathscr{B}}
\newcommand{\sX}[0]{\mathscr{X}}
\newcommand{\sS}[0]{\mathscr{S}}
\newcommand{\sT}[0]{\mathscr{T}}
\newcommand{\cB}[0]{\mathcal{B}}
\newcommand{\cA}[0]{\mathcal{A}}
\newcommand{\cC}[0]{\mathcal{C}}
\newcommand{\cO}[0]{\mathcal{O}}
\newcommand{\cT}[0]{\mathcal{T}}
\newcommand{\cF}[0]{\mathcal{F}}
\newcommand{\cG}[0]{\mathcal{G}}
\newcommand{\ul}[1]{\underline{#1}}
\newcommand{\nat}[0]{\mathrm{nat}}
\newcommand{\tand}[0]{\text{ and }}
\newcommand{\tor}[0]{\text{ or }}
\newcommand{\dcirc}[0]{{\circ\circ}}
\newcommand{\otherwise}[0]{\text{ otherwise }}
\newcommand{\Rat}[0]{\mathfrak{Rat}}
\newcommand{\polring}[1]{[X_1,\cdots,X_{#1}]}
\newcommand{\compring}[1]{\langle X_1,\cdots,X_{#1}\rangle}
\newcommand{\IH}[0]{\mathrm{IH}}
\newcommand{\IC}[0]{\mathrm{IC}}
\newcommand{\PreShv}[0]{\mathrm{PreShv}}
\newcommand{\Shv}[0]{\mathrm{Shv}}
\newcommand{\Perv}[0]{\mathrm{Perv}}
\newcommand{\pt}[0]{\mathrm{pt}}
\newcommand{\fib}[1]{%
  \mathbin{\mathop{\times}\limits_{#1}}%
}
\newcommand{\tens}[1]{%
  \mathbin{\mathop{\otimes}\displaylimits_{#1}}%
}
\newcommand{\comppol}[1]{\langle X_1,\cdots, X_{#1}\rangle}

\title{Perverse Sheaves}
\author{Based on Lectures given by\\\\Prof. Dr. Jens Franke}
\date{SoSe 26, University of Bonn}


\begin{document}
\maketitle
\section{Introduction}
\subsection{Overview of the Course}
%10/12
Given a topological space $X$, we can compute its singular cohomology $H^\bullet(X;k)$ over a commutative ring $k$.
\begin{example}
    \begin{equation*}
        H^i(\{\mathrm{pt}\};k)=\begin{cases}
            k,\quad(i=0),\\
            0,\quad(\otherwise).
        \end{cases}
    \end{equation*}
    \label{example_1_1}
\end{example}
\begin{example}
    \begin{equation*}
        H^i(S^1;k)=\begin{cases}
            k,\quad(i=0,1),\\
            0,\quad(\otherwise).
        \end{cases}
    \end{equation*}
    \label{example_1_2}
\end{example}
\begin{definition}
    \label{definition_1_1}
    Let $X$ be a topological space. The $i$-th Betti number of $X$ is 
    \begin{equation*}
        b_i\defeq \dim_\Q H^i(X;\Q).
    \end{equation*}
\end{definition}
\begin{example}
    \label{example_1_3}
    In Example \ref{example_1_2}, the Betti numbers are 
    \begin{equation*}
        b_i(S^1) = \begin{cases}
            1,\quad(i=0,1),\\
            0,\quad(\otherwise).
        \end{cases}
    \end{equation*}
\end{example}
\begin{theorem}[Poincaré Duality]
    \label{theorem_1_1}
    Let $M$ be a compact oriented manifold without boundary of dimension $d$, then for any $0\leq i\leq d$, we have 
    \begin{equation*}
        H^i(M;\Q)\cong H^{d-i}(M;\Q)^\vee.
    \end{equation*}
    In terms of Betti numbers, this gives us
    \begin{equation*}
        b_i(M) = b_{d-i}(M).
    \end{equation*}
\end{theorem}
\begin{theorem}[Hard Lefschetz]
    \label{theorem_1_2}
    $X$ a smooth complex projective algebraic variety of dimension $\dim_\C X= d$. Then there is an isomorphism 
    \begin{equation*}
    H^i(X,\Q)\to H^{2d-i}(X,\Q)
    \end{equation*}
    given by multiplication with some power of $l\in H^2(X,\Q)$.
    \begin{center}
    \begin{tikzcd}
{H^i(X,\Q)} \arrow[r, "l"] \arrow[rrr, "\sim", bend right] & {H^{i+2}(X,\Q)} \arrow[r, "l"] & \cdots \arrow[r, "l"] & {H^{2d-i}(X,\Q)}
\end{tikzcd}
\end{center}
\end{theorem}
We then get $b_i\leq b_{i+2}$ if $i<d$ and $b_i\geq b_{i+2}$ if $i\leq d$.
\begin{align*}
    b_0\leq b_2\leq\cdots\leq &b_{\mathrm{middle}}\geq\cdots\geq b_{2d},\\
    b_1\leq b_3\leq\cdots\leq &b_{\mathrm{middle}'}\geq\cdots\geq b_{2d-1}.
\end{align*}
\begin{definition}
    \label{definition_1_2}
    For a complex projective variety $X$, we define
    \begin{equation*}
        H^{ij}(X)\defeq H^i(X,\Omega^j_X).
    \end{equation*}
\end{definition}
\begin{theorem}
    \label{theorem_1_3}
    \begin{equation*}
        H^k(X,\C)= \bigoplus_{i+j=k} H^{i,j}(X), \: H^{i,j}\cong \overline{H^{j,i}(X)}.
    \end{equation*}
    Furthermore, 
    \begin{equation*}
        b_k = \sum_{i+k}b_{i,j}, b_{i,j} = b_{j,i}.
    \end{equation*}
\end{theorem}
\begin{corollary}
    In the situation of Theorem \ref{theorem_1_4}, we have $k$ is odd $\Rightarrow$ $b_k$ is even.
\end{corollary}
\begin{notation}
For a topologica space $X$, the constant sheaf of an abelian group $A$ is denoted by $\underline{A}_X$.
\end{notation}
\begin{theorem}[Deligne]
    For be a smooth morphism $f:X\to Y$ between smooth projective complex varieties, we have
    \begin{equation*}
        Rf_*\underline{\Q}_X \cong \bigoplus_{i\in\Z}\mathcal{H}^i(Rf_*\underline{\Q}_X)[-i].
    \end{equation*}
    \label{theorem_1_4}.
\end{theorem}
\begin{remark}
    Consider,
    \begin{equation*}
        S^1\to S^3\stackrel{\substack{ \text{\makebox[0pt][c]{$f$ a Hopf map}}\\\:}}{\to} S^2
    \end{equation*}
    Then $Rf_*\underline{\Q}_{S^3}$. However, $\mathcal{H}^0(Rf_*\underline{\Q}_{S^3}) = \underline{\Q}_{S^2}$ and $\mathcal{H}^1(Rf_*\underline{\Q}_{S^3}) = \Q_{S^2}$. Therefore,
    \begin{equation*}
        Rf_*\underline{\Q}_{S^3}\not\cong \underline{\Q}_{S^2}\oplus\underline{\Q}_{S^2}[-1], H^*(S^3)\not\cong H^*(S^2)\oplus H^{*-1}(S^2).
    \end{equation*}
    Hence $f$ does not satisfy Theorem \ref{theorem_1_4}.
\end{remark}
\begin{question}
Over which space, does the Poincaré duality hold?
\end{question}
This question led to the invention of the following notion.
\begin{definition}[Intersection Cohomology]
    \label{definition_1_3}
    For a variety $X$ over $\C$, we associate an mathematical object $\IH^\bullet(X;\Q)$ which is called intersection cohomology.
\end{definition}

\begin{theorem}
    \label{theorem_1_5}
    For a proper variety $X$ over $\C$ of finite dimension, we have
    \begin{equation*}
        \IH^i(X;\Q)\cong \IH^{2d-i}(X;\Q)^\vee.
    \end{equation*}
\end{theorem} 

Theorem \ref{theorem_1_5} proves some versions of Theorem \ref{theorem_1_2}.
\begin{notation}
    Let $X$ be a topological space. The derived category of constructible sheaves on $X$ is denoted by $D^b_c(X)$.
\end{notation}
\begin{theorem}
    \label{theorem_1_6}
    For a proper variety $X$ over $\C$ of finite dimension, there is a complex of sheaves $\IC_X\in D^b_c(X)$ such that% is a perverse sheaf $\IC_X\in D^b_c(X)$ in the derived category of constructible sheaves on $X$ such that 
\begin{equation*}
    \IH^*(X,\Q) \cong H^*(X,\IC_X).
\end{equation*}
\end{theorem}
\begin{definition}
    Such $\IC_X$ in Theorem \ref{theorem_1_6} is called a perverse sheaf.
\end{definition}
%which is the cohomology of this complex of sheaves.
%Perverse sheaves, Belinson-Bernstein-Deligne 1983, BBD + Gaber also extended Deligne's theorem to $f:X\to Y$ proper.
\begin{remark}
    Perverse sheaves were introduced by Belinson-Bernstein-Deligne in 1983. Together with Gabber, they extended Theorem \ref{theorem_1_4} to the case when $f:X\to Y$ is proper.
\end{remark}
\subsection{Preliminaries}

\begin{notation}
    Let $X$ be a topological space and $k$ be a commutative ring. The category of presheaves of $k$-modules over $X$ is denoted by $\PreShv(X;k)$, similarly the sheaves of $k$-modules over $X$ is denoted by $\Shv(X;k)$.
\end{notation}
\begin{example}$\:$
        For all $M\in\mathrm{Mod}(k)$, there is a constant sheaf $\underline{M}_X\in\Shv(X;k)$.
\end{example}
\begin{definition}
    A local system on a topological space $X$ is a sheaf on $X$ which is locally a constant sheaf of an abelian group.
\end{definition}
\begin{example}
    % For $r>0$, consider morphisms of groups,
    %     \begin{equation*}
    %         p:\pi_1(X,x_0)\to \Aut(k^{\oplus r}).
    %     \end{equation*}
    %     Define $L_p4$ as follows. 
    For $r>0$, consider morphisms of groups,
        \begin{equation*}
            p:\pi_1(X,x_0)\to \Aut(k^{\oplus r}).
        \end{equation*}
        Define $L_p$ as follows. Let $\pi:\tilde{X}\to X$ be the universal covering space. For $U\subseteq X$, 
        \begin{equation*}
            L_p \defeq \{\pi^{-1}(U)\to k\sep \forall g\in\pi_1(X,x_0), f(x\cdot g) = p(g)^{-1}f(x)\}.
        \end{equation*}
        Then $L_p$ is a local system on $X$.
\end{example}
\begin{theorem}
    $\Shv(X;k)$ is an abelian category.
\end{theorem}

\begin{proposition}
    The canonical inclusion $\Shv(X;k)\to \PreShv(X;k)$ admits a left adjoint and such functor is called the sheafification,
    \begin{equation*}
    \PreShv(X;k)\ni \cF\mapsto\cF^\dagger\in\Shv(X;k).
    \end{equation*}
    Explicitly, for $\cF\in\PreShv(X;k)$ and $\cG\in\Shv(X;k)$, we have 
    \begin{equation*}
        \mathrm{Mor}_{\PreShv}(\cF,\cG)\cong \mathrm{Mor}_{\Shv}(\cF^\dagger,\cG),
    \end{equation*}
    which is natural. In particular, setting $\cG = \cF^\dagger$, there is a natural map $\cF\to\cF^\dagger$ which is obtained from $\id\in\mathrm{Mor}_{\Shv}(\cF^\dagger,\cF^\dagger)$.
\end{proposition}
\begin{proposition}
    The natural morphism $\cF\to \cF^\dagger$ induced isomorphisms on stalks.
\end{proposition}

\begin{definition}
        Given a complex of sheaves,
        \begin{equation*}
            \cF^\bullet \defeq \cdots\to \cF^{i-1}\stackrel{d^{i-1}}{\to}\cF^i\stackrel{d^{i}}{\to}\cF^{i+1}\cdot
        \end{equation*}
        we define its cohomology of sheaves to be 
        \begin{equation*}
            \cH(\cF^\bullet) \defeq\ker d^i\oplus\image d^{i-1}.
        \end{equation*}
        A morphism of complexes is called a quasi-isomorphism if it induces isomorphisms in cohomology.
\end{definition}

\begin{definition}
    The derived category on $\Shv(X;k)$ is a category satisfying the following.
    \begin{enumerate}[i. ]
        \item The objects are complex of sheaves,
        \begin{equation*}
            \cF^\bullet \defeq \cdots\to \cF^{i-1}\stackrel{d^{i-1}}{\to}\cF^i\stackrel{d^{i}}{\to}\cF^{i+1}\cdot
        \end{equation*}
        \item The morphisms are morphisms of complexes.
        \item Quasi-isomorphisms are invertible.
    \end{enumerate}
\end{definition}

\begin{remark}
    The derived category of sheaves on $X$ over $k$ is a triangulated category.
\end{remark}

However the derived category turns out to be to large to handle. Therefore, we will use a smaller category called the bounded constructible derived category of sheaves denoted by $D_c^b(X;k)$.

\begin{remark}
    The category of bounded constrictible derived category of sheaves $D_c^b(X)$ is a subcategory of the derived category of sheaves which is more manageable.
\end{remark}
\begin{definition}
    There exists a stratification in locally closed subspaces $X = \bigsqcup_\alpha S_\alpha$. That is for $\cF^\bullet\in D_c^b(X)$, we have $\cH^i(\cF^\bullet)|_{S_\alpha}$ is a finite rank local system and $\cH^i(\cF^\bullet) = 0$ for $\ab{i}\gg 0$.
\end{definition}

\begin{remark}
    Let $f:X\to Y$ be a continuous map. The following so-called six-functors are important when it comes to produce example of sheaves,
    \begin{equation*}
        Rf_*\dashv f^*, f_!\dashv f^!, (\cdot\otimes\cdot)\dashv \Hom(\cdot,\cdot).
    \end{equation*}
\end{remark}

% \begin{center}
% \begin{tabular}{l|l}
% Left Adjoints         & Right Adjoints              \\ \hline
% $Rf_*$ pushforward    & $f^*$ pullback              \\ \hline
% $f_!$                 & $f^!$                       \\ \hline
% $(\cdot\otimes\cdot)$ & $\mathrm{Hom}(\cdot,\cdot)$
% \end{tabular}
% \end{center}

\begin{remark}
    \label{remark_1_6}
    If $f$ is proper then $Rf_* = f_!$, if $f$ is open immersion then $f^* = f^!$.
\end{remark}

%10/13

\begin{theorem}
    \begin{equation*}
        D_c^b(X;k)\cong D^b(\Shv_c(X;k)).
    \end{equation*}
\end{theorem}

\begin{example}
    For a scheme $X$ and $\alpha:X\to\mathrm{pt}$, we have,
    \begin{equation*}
        \alpha_*\ul{k}_{\mathrm{pt}} \cong \ul{k}_X,\alpha^!\ul{k}_{\mathrm{pt}}\cong \omega_X.
    \end{equation*}
\end{example}

With the aid of these functors, we define(reconstruct) following complexes.
\begin{definition}
    \begin{align*}
        H^\bullet(X;k) &\defeq \alpha_*\alpha^*\ul{k}_\pt,\\
        H^{\mathrm{BM}}(X;k) &\defeq\alpha_*\alpha^!\ul{k}_\pt,\\
        H^*_c(X;k)&\defeq \alpha_!\alpha^*\ul{k}_\pt,\\
        H_*(X;k)&\defeq \alpha_!\alpha^!\ul{k}_\pt.
    \end{align*}
\end{definition}

\begin{example}
    Let $Z\subseteq X$ be closed and $U=X\backslash Z$. Consider $Z\stackrel{i \text{ closed}}{\hookrightarrow}X\stackrel{j \text{open}}{\hookleftarrow}U$. This gives natural transformations,
    \begin{equation*}
        (\cdot)\to i_*i^*(\cdot), (\cdot)\to j_*j^*(\cdot),(\cdot)\to i_!i^!(\cdot),(\cdot)\to j_!j^!(\cdot).
    \end{equation*} 
    Furthermore, this gives an exact triangle,
    \begin{equation*}
        j_!j^!(\cdot)\to (\cdot)\to i_*i^*(\cdot)\to[+1].
    \end{equation*}
    Using Remark \ref{remark_1_6} and plugging in $\ul{k}_X = \alpha_*\ul{k}_\pt$, we obtain
    \begin{equation*}
        \alpha_*j_!\ul{k}_U\to\alpha_*\ul{k}_X\to \alpha_*i_*\ul{k}_Z\to
    \end{equation*}
\end{example}

Suppose $X=U\cup V$ and $W = U\cap V$, and there are natural transformations,
\begin{equation*}
    (\cdot)\to (j_V)_*(j_U)^*(\cdot).
\end{equation*}

In the category of sheaves $\Shv(X;k)\subseteq D_c^b(X;k)$.

\begin{center}
    \begin{tikzcd}
\Shv(X;k) \arrow[r, hook]                   & D^b(\Shv(X;k))                   \\
\Shv_c(X;k) \arrow[u, hook] \arrow[r, hook] & D^b_c(\Shv(X;k)) \arrow[u, hook]
\end{tikzcd}
\end{center}

\begin{definition}
    The Perverse sheves is $\Perv(X;k)\subseteq D_c^b(\Shv(X;k))$, where $X$ is fnite-type variety over a field.
    \par Explicitly, for $D(\cdot)= \mathcal{H}om(\cdot,\omega_X)$, define $\cF\in D_c^b(X;k)$ is a perverse sheaf 
    if for all $n\in\Z$, $\dim\supp\mathcal{H}(\cF)\leq i$ and $\dim\supp\cH(D(\cF))\leq i$.
\end{definition}

\begin{example}
    Let $X$ be a smooth variety over $\C$ of dimension $\dim_\C X = d$, then,
    \begin{equation*}
        \omega_X = k_X[2d],
    \end{equation*}
    then $\underline{k}_X[d]$ is a perverse sheaf and $\mathbb{D}(\underline{k}_k[d]) = \underline{k}_X[d]$.
    As,
    \begin{equation*}
        \cH(k_x[d]) = \begin{cases}
            \underline{k}_X \quad (i=d),\\
            0,(\otherwise).
        \end{cases}
    \end{equation*}
    Hence 
    \begin{equation*}
        \dim\supp(\cH(\underline{k}_X[d]))=\begin{cases}
            d,\quad(i=d),\\
            -\infty,\quad(\otherwise).
        \end{cases}
    \end{equation*}
    Similarly, if $L $is a local system on $X$, then $L[d]\in\Perv(X;k)$.
\end{example}

\begin{theorem}
    $\Perv(X;k)$ is an Abelian category which is both Noetherian and Artinian.
\end{theorem}

Every perverse sheaf has a composition series with ?
All simple objects are constructed as follows,

\begin{equation*}
    \underbrace{X}_{\text{smooth finite-type variety}}\hookleftarrow \underbrace{Z}_{\text{closed subvariety}}\hookleftarrow \underbrace{U}_{\text{open smooth}}.
\end{equation*}

Let $L$ be a local system on $V$, then we construct a perverse sheave,

\begin{align*}
    \IC_Z(L)\subseteq\Perv(X;k),
\end{align*}

If it is supported on $Z$ and 
\begin{equation*}
    \IC_Z(L)|_V = L[\dim V].
\end{equation*}

\begin{thebibliography}{9}
\bibitem{Achar}
Pramod Achar, Perverse Sheaves and Applications to Representation Theory.
\bibitem{Cataldo}
Mark de Cataldo, Luca Migliorini, The Decomposition Theorem and the topology of algebraic maps.
\end{thebibliography}
\end{document}